% Review-language discussion and limitations (M.N.).
Comparing the network scores with review language answers our main question: home cooks do not call a recipe authentic simply because it lists typical single ingredients ($\rho = 0.019$, $z = 0.57$), but because it combines ingredients in the cuisine's characteristic way ($\rho = 0.091$, $z = 6.80$, $p = 0.0005$). The signal is positive in all nine cuisines, survives holding review volume fixed at one, and is confirmed by a rule-free sentence-embedding classifier. Yet the modest magnitude ($\rho \approx 0.09$) shows that ingredient lists alone cannot capture everything ``authentic'' means: identical ingredients yield different dishes depending on cooking technique and proportions, and Food.com's predominantly North American, English-speaking reviewers judge authenticity through personal memories and familiar diaspora or restaurant versions rather than local culinary tradition.

Four NLP limitations qualify this result. First, the Word2Vec screen introduces upstream sensitivity to seed choice, neighbor cutoff, and curation, and misses purely implicit endorsements with no lexicon term (e.g., ``just like my grandmother made''). Second, within the screened pool, four explicit KWIC rules reach 84.2\% precision and 72.7\% recall on \texttt{test\_100}, missing implicit praise and occasionally flagging branded ingredients or references to other recipes. Third, explicit denials (``not authentic'') are grouped with neutral mentions rather than modeled as a separate negative class. Finally, validation relied on one annotator; double-coding for inter-annotator agreement and auditing implicit judgments with an instruction-tuned language model are natural next steps.
